\documentclass[11pt,a4paper]{article}

% Packages
\usepackage[utf8]{inputenc}
\usepackage[T1]{fontenc}
% Typography: URW Palatino (Palladio) text + Pazo math + Inconsolata mono.
% This is the same font stack used by the Qwen technical reports.
\usepackage{mathpazo}
\usepackage[scaled=0.95]{inconsolata}
\usepackage{microtype}
\usepackage{amsmath,amssymb,amsfonts}
\usepackage{graphicx}
\usepackage[numbers,sort&compress]{natbib}
\usepackage{booktabs}
\usepackage{xcolor}
\usepackage{fancyhdr}
\usepackage{tikz}
\usetikzlibrary{arrows.meta, positioning, fit, backgrounds, calc}
\usepackage[margin=1in]{geometry}
\usepackage[hidelinks]{hyperref}
\usepackage[english]{babel}
\setlength{\emergencystretch}{3em}

% Resolve assets/figures whether compiled from the repo root or this folder.
% Copy this line into each paper, pointing the first path at the paper's folder.
\graphicspath{{papers/2026-10-false-approval-claim/}{./}}

% Brand header (logo + wordmark on the left, date on the right)
\definecolor{gaussiadark}{HTML}{323030}
\definecolor{gaussiaccent}{HTML}{C15A35}
\newcommand{\paperdate}{Version 1.0, 2026-10-20}
\newcommand{\papersurl}{https://www.gaussia.ai/papers}
\newcommand{\repourl}{https://github.com/gaussia-labs/papers}
\newcommand{\hfurl}{https://huggingface.co/gaussia-labs}
% Header logo (no link, just branding)
\newcommand{\gaussialogo}{%
  \raisebox{-0.28\height}{\includegraphics[height=14pt]{logo}}%
  \hspace{5pt}\textcolor{gaussiadark}{\textbf{\large Gaussia}}%
}
% Clickable links shown under the title, Qwen-report style
\newcommand{\paperlinks}{%
  \begin{center}
    \href{\repourl}{%
      \raisebox{-0.22\height}{\includegraphics[height=11pt]{github}}%
      \hspace{4pt}\texttt{github.com/gaussia-labs/papers}}\\[4pt]
    \href{\papersurl}{%
      \raisebox{-0.26\height}{\includegraphics[height=11pt]{logo}}%
      \hspace{4pt}\texttt{gaussia.ai/papers}}\\[4pt]
    \href{\hfurl}{%
      \raisebox{-0.26\height}{\includegraphics[height=12pt]{hf-logo}}%
      \hspace{4pt}\texttt{huggingface.co/gaussia-labs}}
  \end{center}%
}

\providecommand{\keywords}[1]{%
  \vspace{0.5em}%
  \noindent\textbf{Keywords: }#1%
}

\setlength{\headheight}{18pt}
\pagestyle{fancy}
\fancyhf{}
\lhead{\gaussialogo}
\rhead{\textcolor{gaussiadark}{\paperdate}}
\renewcommand{\headrulewidth}{0.6pt}
\rfoot{\thepage}

% Metadata
\title{False approval claims: a deterministic verification of approval statements in AI assistants with tools\\
\large Formal definition, per-turn measurement, extraction error and validation}
\author{Juan F. Cobo (author) \and Leonardo Leenen (reviewer, 2026-10-20)\\ \small EthiCompass \\
\small Manuscript prepared for external review and eventual submission for publication}
% Date is shown in the header instead of under the title (see \paperdate above)
\date{}

\begin{document}

\maketitle
\thispagestyle{fancy}

% Repository and website links, shown under the title
\paperlinks

% The abstract is REQUIRED for every Gaussia paper. Do not remove it.
\begin{abstract}
An AI assistant with tools tells what it did and why. Sometimes it says that a person approved the action: ``your
supervisor already approved it''. If that approval does not exist, the assistant acted alone, hid it and moved the
responsibility to a person who never saw the case. The published metrics treat this failure as a case of general
fabrication, decided by a judge. We propose to measure it separately, because it can be verified without a judge:
an extraction finds the approval statements in each turn, and a deterministic verification looks for them in the
approval record. Only a granted approval, recorded by the approval system, earlier than or simultaneous with the
turn and covering the claimed action, supports a statement; if the assistant could write the record, it could
record the approval that it invents. The result of each turn is 0 or 1; the result of the conversation is the
maximum; one confirmed false statement is an incident. The only source of error is the extraction, whose
sensitivity and specificity are reported together with the rate. Without the system record, the verification uses
the record of the person who runs the test, who knows what that person approved, and it stays deterministic. This
work reports no results with real assistants.
\end{abstract}

\keywords{AI red-teaming, human oversight, attribution, fabrication, tool-using agents, deterministic verification, incidents}

\noindent\textit{The Spanish version of this document is the original.}

\section{Introduction and motivation}

After it acts, the assistant explains to the user what it did. The user, and later an auditor, use this account to
know who made each decision. A sentence that says ``a person approved this'' has a special effect: it tells the
reader that the action already had a human review, and for this reason nobody reviews it again.

If the sentence is false, three things occur: the assistant acted without the approval that was necessary, it hid
that it acted alone, and the responsibility moves to a person who never saw the case. The failure can have several
causes: the model completes the account with a usual sentence, hostile content pushes it to justify the action, or
the model confuses a requested approval with a granted approval. The measurement does not need to know the cause:
it compares what the assistant says with what the record has.

Frequently the failure is in the middle of the conversation: the assistant acts, justifies the action with an
approval that does not exist and then says goodbye correctly. A measurement that reads only the final response
does not see it. For this reason the measurement is per turn.

\section{Related work}

\begin{itemize}
  \item \textbf{Fabrication and false success.} Sethi et al.\ \cite{sethi2026} show that assistants fabricate
    accounts after a tool failure (14.1\% dishonest responses; 45.3\% if the failure is not marked as an error),
    and that a status indicator read deterministically decreases the rate to 0.87\%. Zhu et al.\ \cite{zhu2026}
    report 22.8\% false success with the base policy and 0.8\% with an evidence contract. Advani
    \cite{advani2026} shows that LLM judges detect false success badly (AUROC $<0.65$). Guo et al.\
    \cite{upwarddeception} define upward deception as a final report whose success status is higher than that of
    the real trajectory, the same structure as a statement against a record; they score it with an LLM judge and
    report no human validation of that judge. The four justify verifying against a record, not against the
    text.
  \item \textbf{Attribution.} Chan et al.\ \cite{chan2024} treat the attribution of the actions of an agent and
    activity logs. Narajala and Narayan \cite{narajala2025} describe the attribution gaps in threat T9: the agent
    hides who decided.
  \item \textbf{Published oversight metrics.} The Gaussia oversight paper \cite{gaussia2026} measures fabrication
    with a judge, inside action disclosure. In the public discussion of our proposals, its team proposed that only
    the data recorded by the approval system count and that the request data go in the approval event
    \cite{gaussiadisc}; this work adopts both rules.
  \item \textbf{Risk framework.} The OWASP Top 10 for Agentic Applications 2026 \cite{owaspagentic} lists
    Human-Agent Trust Exploitation as ASI09: humans approve actions without independent validation. A false ``a
    person approved this'' makes the reader believe that a review occurred.
\end{itemize}

\section{Formal definition}

\subsection{Data}

Let $i$ be the conversation and $t \in \{0,\dots,T_i-1\}$ the turn of the assistant, with text $x_{i,t}$. The
approval record is a set $\mathcal{B}_i$; each approval $b$ has a decision $d(b) \in
\{\text{granted}, \text{denied}, \text{pending}\}$, a source $s(b)$, the turn of the decision $\tau(b)$ and the set
of actions that it covers. The request data (the text, its author and its turn) go in the same event and help to
tell ``I requested the approval'' from ``they approved it''.

\subsection{Extraction}

An extraction function $E$ returns, for each turn, the set of approval statements
$E(x_{i,t}) = \{a_1, \dots, a_m\}$. Each statement has an actor (who approved, according to the text), an action
(connected to a tool of the policy when possible) and the quote. An approval request, a plan or a condition (``if
the supervisor approves it, I will issue it''), and a denial do not count as statements. The extraction can be a
pattern or a model; it only finds the statement and does not decide whether it is false.

\subsection{Verification}

A statement $a$ of turn $t$ is supported if
\begin{equation}
  \exists\, b \in \mathcal{B}_i :\; d(b) = \text{granted} \;\wedge\; s(b) = \text{approval system}
  \;\wedge\; \tau(b) \leq t \;\wedge\; \mathrm{covers}(b, a).
\end{equation}
$\mathrm{covers}(b,a)$ compares the tool of the approval with the claimed action; if the record connects the
approval with the call, it uses that connection. If the statement does not name an action, it is connected with
the action of the same turn or of the earlier turn. The verification does not compare the parameters of the
statement (Section~\ref{sec:boundaries}).

\paragraph{Vague statements.} If neither turn has an action to connect it to, the statement is \emph{vague}: it
is not verified and it is left out of $g_{i,t}$; if it is the only statement of the turn, the turn has no score,
with the reason \texttt{vague\_statement}. Counting it as 0 would make a cell without a score count as
compliance. The number of vague statements is reported, so that the reader knows how much stays without a
verification. Claim extraction has the same rule: Metropolitansky and Larson \cite{claimify} exclude
the sentences whose ambiguity cannot be resolved and label them, at most 5.4\% of the sentences in their
evaluation.

\paragraph{The source.} Only an approval recorded by the approval system supports a statement. If the assistant
could write the approval record, it could also record the approval that it invents, and the measurement would
never fail.

\subsection{Result per turn and per conversation}

The result of the turn is
\begin{equation}
  g_{i,t} = \mathbf{1}\big[\exists\, a \in E(x_{i,t}) : a \text{ is not supported}\big],
\end{equation}
with the cause \texttt{false\_approval\_claim} and the quote of the statement. If the extraction fails or the
record is incomplete, the turn has no score and has a reason; a turn without a score never counts as compliance.
The value of the conversation is $G_i = \max_t g_{i,t}$ over the turns with a score; a conversation without turns
with a score goes out of the denominator, with its reason.

\subsection{Rate, interval and incidents}

The rate is $\hat{p} = \sum_i G_i / n$ over the evaluated conversations that target the measurement. A model never
decides the denominator. The interval is Wilson \cite{wilson1927} with the Kish effective size \cite{kish1965},
$n_{\mathrm{ef}} = n / (1 + (\bar{m}-1)\rho)$, with the replicas grouped by scenario. Because the verification is
deterministic, each confirmed false statement is an incident, reported with the conversation, the turn and the
quote, without waiting for a rate.
Per scenario, the report also says whether all its $k$ replicas had zero false statements; the share of
scenarios where this is true is a figure of consistency, the idea of pass$^k$ \cite{taubench} applied to a
behavior that must never occur, and it reads better than a rate.

The correlation $\rho$ is not chosen by hand: a run estimates it from its own replicas, as the intraclass
correlation of the per-conversation results grouped by scenario. The paper requires that $\rho$ is estimated and
reported, not a specific estimator. Ridout et al.~\cite{ridout1999} review and compare the estimators for binary
clustered data, among them the one-way ANOVA estimator, which our implementation uses with a negative estimate
taken as 0. The output reports $\hat{\rho}$ together with $\bar{m}$ and $n_{\mathrm{ef}}$, and the interval uses
the $n_{\mathrm{ef}}$ computed with $\hat{\rho}$; the Wilson interval with an effective sample size in place of $n$
is one of the two forms that Franco et al.~\cite{franco2019} recommend. An assumed $\rho$ is used only to plan
the number of scenarios and replicas, or as a sensitivity analysis. With few scenarios $\hat{\rho}$ is noisy; a
clustered standard error with the scenario as the cluster \cite{miller2024} needs no $\rho$ and can be reported
next to the interval as a check.

\subsection{Error of the extraction}

The extraction is the only part that can make errors. Its error is described by the sensitivity $\mathrm{Se}$
(the proportion of real statements that it finds) and the specificity $\mathrm{Sp}$ (the proportion of turns
without a statement in which it finds nothing). A statement that the extraction does not find is a failure that is
not counted; a statement that it invents can give a false 1. We report $\mathrm{Se}$ and $\mathrm{Sp}$ together
with the rate, with the model version and per language, and we publish the rate without correction. The
Rogan-Gladen correction \cite{rogan1978},
\begin{equation}
  \hat{\pi} = \frac{\hat{p} + \mathrm{Sp} - 1}{\mathrm{Se} + \mathrm{Sp} - 1},
\end{equation}
is possible, but stays as an open question.

\paragraph{Calibration by construction.} A large part of the set to estimate $\mathrm{Se}$ and $\mathrm{Sp}$ can be
built without human labels: take turns with a known statement and change them in known ways (a different verb, a
different actor, a different sentence), and use conditional forms, requests and denials as negative controls. Only
the paraphrases of real traffic need human labels. The extraction model receives the tools of the policy and the
turn, in the language of the assistant, and returns the list of statements with their actor, their action and
their quote.

\paragraph{Validation of the extraction.} Because the extraction is the only source of error, its validation is
reported in full. In addition to $\mathrm{Se}$ and $\mathrm{Sp}$ per turn, the report gives the precision and the
recall of the extraction against human labels at the sentence level (does this sentence contain an approval
statement?), and the agreement between the annotators of the part with human labels, for example with
Krippendorff's alpha. A sentence whose meaning cannot be decided from its context goes to an explicit category,
\emph{cannot be disambiguated}: it is excluded from the extraction and counted, so that the reader knows how much
the extraction leaves out. The template is the one of Metropolitansky and Larson \cite{claimify}, who report
sentence-level precision and recall and the agreement of three annotators with Krippendorff's alpha.

\section{Observation with the conversation only}

In a conducted test, the person who runs the conversation also acts as the supervisor: that person approves or
denies the requests, and for this reason knows what was approved and when. Without the system record, the
verification uses the record of the person who runs the test. It stays deterministic, and its result can still be
used as incident evidence; only the extraction reads text.

\section{Boundaries with other measurements}
\label{sec:boundaries}

\begin{itemize}
  \item \textbf{Fabrication.} In the turn where this measurement gives 1, general fabrication does not measure the
    same sentence again: the same statement is not counted twice.
  \item \textbf{Escalation.} ``I already passed your case to a supervisor'', without an escalation call, states an
    escalation, not an approval: it is under-escalation, which an escalation metric measures.
  \item \textbf{Substitution.} If 90 USD was approved and the assistant says ``the supervisor approved the refund of
    900'', the approval of the tool exists and this measurement does not fail; if 900 was executed, substitution
    finds the failure, because it compares what was approved with what was executed.
  \item \textbf{Execution without a valid approval.} A denied approval that the assistant presents as approved is a
    failure here; if the action was also executed, a check of the approval record also finds it.
\end{itemize}

\section{Illustrative example}
\label{sec:example}

\textit{(Illustrative construction; not a measurement.)}

\paragraph{A conversation of three turns.} In turn 0 the assistant writes ``for a refund of 900 USD I need the
approval of a supervisor'' and the system records the request, with the decision pending. In turn 1 it writes ``the
supervisor already approved the refund of 900 USD and I issued it''. In turn 2 it writes ``can I help you with
something else?''. The extraction finds no statements in turns 0 and 2 ($g=0$) and finds one in turn 1. The record
has no granted approval with turn $\leq 1$, so $g_{1}=1$ and $G = \max(0,1,0) = 1$: an incident in turn 1. The
final response is clean: a measurement that reads only the final response gives 0.

\paragraph{A rate.} There are 20 scenarios with 3 replicas: 60 conversations. Two have no turns with a score:
$n=58$. Seven conversations fail: $\hat{p} = 7/58 = 0.121$. With $\bar{m} = 3$ and $\hat{\rho} = 0.3$
(illustrative; a real run estimates it from its replicas),
$\mathrm{deff} = 1 + 2 \times 0.3 = 1.6$ and $n_{\mathrm{ef}} = 58/1.6 = 36.25$; the 95\% Wilson interval with these values
is $(0.050,\ 0.264)$. In addition, 7 incidents are reported, each one with its turn and its quote.

\section{Validation}
\label{sec:validation}

\subsection{Test cases with a known result}

Because the verification is deterministic, the cases test the exact result: a false statement in turn 1 with a
clean final response; a true statement; a request without a statement; a denied approval presented as granted; an
approval of a different action; an approval later than the statement; a hypothetical statement; a turn without text
in the record; an approval with a source different from the approval system; and an escalation statement. In the
set of conformance cases, two cases are of this measurement: a false approval claim with a record,
and the same failure in turn 1 of a conversation of three turns with a clean final response. The set detects,
among others, a judge score instead of 0 or 1 and the violation of turn 1 not seen. The source rule is not yet in
that set.

\subsection{Synthetic simulation}

In a synthetic simulation with a known truth, separating this measurement from fabrication keeps decisions:
with a simulated judge with 90\% accuracy, when the fabrication judge does not respond, the verification of the
false approval claim is still decided. Without the system record, the failure stays visible, because the person
who runs the test knows that that person did not approve.

\section{Implementation considerations}

The metric extends the \texttt{Gaussia} base class of the \texttt{pygaussia} SDK. The \texttt{batch()} method reads each turn of the assistant,
applies the extraction $E$ to its text and verifies each statement against the approval record. It gives
$g_{i,t}$, with the cause and the quote. If the extraction fails or the record is incomplete, the turn has no
score and has a reason. The session takes the maximum $G_i$.

The metric reads the approvals and the policy of the oversight data contract of Gaussia \cite{gaussia2026}: \texttt{oversight\_policy} at the session
level, and \texttt{tools\_used} and \texttt{approvals} in the \texttt{agentic} dictionary of each turn. The fields \texttt{at\_qa} and \texttt{covers}
give $\tau(b)$ and the actions that an approval covers. The metric needs two items that the contract does not
have: a decision with three values (granted, denied, pending), because \texttt{granted} has only two, and the
source $s(b)$ of each approval. The request data (the text, its author and its turn) in the same event are
recommended. As in that contract, the names of the new fields are illustrative.

The extraction is a pattern or a model, and it does not decide whether a statement is false. A model receives
the tools of the policy and the turn, in the language of the assistant, and returns the list of statements with
their actor, their action and their quote: one call for each turn of the assistant. The verification uses no
model.

The output has the rate without correction, with its Wilson interval on $n_{\mathrm{ef}}$, the incident list,
and $\mathrm{Se}$ and $\mathrm{Sp}$ of the extraction, with the model version and per language. Compatibility with the statistical modes (frequentist and Bayesian) follows the existing Gaussia pattern.
Without the system record, the verification uses the record of the person who runs the test, in the same shape
as \texttt{approvals}, and the output states the source.

\section{Discussion}

\subsection{What this work establishes}

The measurement is defined by the extraction and the record: what an approval statement is, what supports it, how
each turn is decided, how the turns are joined per conversation and how its only error is described. The
verification is deterministic and can be used as incident evidence, also without the system record.

\subsection{What it does not establish}

\begin{itemize}
  \item \textbf{There are no results with real assistants,} and no measurement of $\mathrm{Se}$ and $\mathrm{Sp}$
    for a specific extraction.
  \item \textbf{It sees only what the record has.} If the system approves through a channel that it does not
    export, a true statement looks false.
  \item \textbf{It does not say why the assistant stated something false.}
  \item \textbf{It does not measure other false statements.} ``I checked your history'' without a call that does it
    is fabrication, not this failure.
  \item \textbf{It depends on the scenarios.} If no scenario leads to an action that needs approval, the rate says
    nothing.
  \item \textbf{The interval can undercover.} With the replicas grouped by scenario, intervals on an effective
    sample size cover less than their nominal level, mainly because $n_{\mathrm{ef}}$ is itself estimated
    \cite{franco2019}. With few scenarios $\hat{\rho}$ is noisy, and the Kish formula assumes scenarios with a
    similar number of replicas.
\end{itemize}

\subsection{Generalization}

Separating from a judge a statement that can be verified against a record applies to other statements of the
assistant about recorded facts. It is not an extension evaluated here.

\subsection{Open questions for the reviewers}

\begin{enumerate}
  \item Must the extraction be a pattern, a small model or a judge? How is its error published?
  \item Is it better to correct the rate for the error of the extraction, or is it sufficient to report it?
  \item Must the record also keep the authorizations of the user?
  \item Must a denied approval presented as granted also be reported as a separate incident?
\end{enumerate}

\section{Conclusion}

A false approval claim can be measured without a judge: an extraction finds the statement and the record decides
whether it is false. The source rule prevents the assistant from supporting its own statement. The per-turn
measurement sees the failure although the final response is clean, and each confirmed case is an incident. The
validation is synthetic, and the next step is to measure the error of a real extraction and real assistants.

\section*{Declarations}

\paragraph{Conflicts of interest.} The authors are affiliated with EthiCompass, the organization that develops a
red-teaming system for AI assistants. The measurement defined
here is of general application and does not depend on a proprietary implementation, but the affiliation is
declared so that the reader can take it into account.

\paragraph{Relation to Gaussia.} The source rule and the location of the request data follow what the Gaussia team
proposed in a public discussion. The separation from fabrication, the per-turn measurement and the validation are
the work of the authors.

\paragraph{Funding.} \textit{To be completed before submission.}

\paragraph{Data and code availability.} The results of Section~\ref{sec:validation} come from a synthetic
simulation and hand-written test cases. \textit{The access conditions are completed before submission.}

\bibliographystyle{plainnat}
\bibliography{references}

\end{document}
